% Options for packages loaded elsewhere
\PassOptionsToPackage{unicode}{hyperref}
\PassOptionsToPackage{hyphens}{url}
\PassOptionsToPackage{dvipsnames,svgnames,x11names}{xcolor}
\documentclass[
  10pt,
]{article}
\usepackage{xcolor}
\usepackage[margin=2.2cm]{geometry}
\usepackage{amsmath,amssymb}
\setcounter{secnumdepth}{-\maxdimen} % remove section numbering
\usepackage{iftex}
\ifPDFTeX
  \usepackage[T1]{fontenc}
  \usepackage[utf8]{inputenc}
  \usepackage{textcomp} % provide euro and other symbols
\else % if luatex or xetex
  \usepackage{unicode-math} % this also loads fontspec
  \defaultfontfeatures{Scale=MatchLowercase}
  \defaultfontfeatures[\rmfamily]{Ligatures=TeX,Scale=1}
\fi
\usepackage{lmodern}
\ifPDFTeX\else
  % xetex/luatex font selection
\fi
% Use upquote if available, for straight quotes in verbatim environments
\IfFileExists{upquote.sty}{\usepackage{upquote}}{}
\IfFileExists{microtype.sty}{% use microtype if available
  \usepackage[]{microtype}
  \UseMicrotypeSet[protrusion]{basicmath} % disable protrusion for tt fonts
}{}
\makeatletter
\@ifundefined{KOMAClassName}{% if non-KOMA class
  \IfFileExists{parskip.sty}{%
    \usepackage{parskip}
  }{% else
    \setlength{\parindent}{0pt}
    \setlength{\parskip}{6pt plus 2pt minus 1pt}}
}{% if KOMA class
  \KOMAoptions{parskip=half}}
\makeatother
\usepackage{longtable,booktabs,array}
\usepackage{caption}
\captionsetup[table]{skip=6pt}
\newcounter{none} % for unnumbered tables
\usepackage{calc} % for calculating minipage widths
% Correct order of tables after \paragraph or \subparagraph
\usepackage{etoolbox}
\makeatletter
\patchcmd\longtable{\par}{\if@noskipsec\mbox{}\fi\par}{}{}
\makeatother
\usepackage{graphicx}
\makeatletter
\newsavebox\pandoc@box
\newcommand*\pandocbounded[1]{% scales image to fit in text height/width
  \sbox\pandoc@box{#1}%
  \Gscale@div\@tempa{\textheight}{\dimexpr\ht\pandoc@box+\dp\pandoc@box\relax}%
  \Gscale@div\@tempb{\linewidth}{\wd\pandoc@box}%
  \ifdim\@tempb\p@<\@tempa\p@\let\@tempa\@tempb\fi% select the smaller of both
  \ifdim\@tempa\p@<\p@\scalebox{\@tempa}{\usebox\pandoc@box}%
  \else\usebox{\pandoc@box}%
  \fi%
}
% Set default figure placement to htbp
\def\fps@figure{htbp}
\makeatother
\setlength{\emergencystretch}{3em} % prevent overfull lines
\providecommand{\tightlist}{%
  \setlength{\itemsep}{0pt}\setlength{\parskip}{0pt}}
\usepackage{fontspec}
\setmainfont{STIX Two Text}
\setmonofont{Menlo}[Scale=0.85]
\usepackage{xeCJK}
\setCJKmainfont{Songti SC}
\usepackage{microtype}
\usepackage{etoolbox}
\AtBeginEnvironment{longtable}{\footnotesize}
\usepackage{xurl}
\setlength{\emergencystretch}{3em}
\sloppy
\setCJKmonofont{Songti SC}
\usepackage{unicode-math}
\setmathfont{STIX Two Math}
\usepackage{newunicodechar}
\newunicodechar{→}{\ensuremath{\rightarrow}}
\newunicodechar{∎}{\ensuremath{\blacksquare}}
\newunicodechar{≠}{\ensuremath{\neq}}
\newunicodechar{≈}{\ensuremath{\approx}}
\newunicodechar{↔}{\ensuremath{\leftrightarrow}}
\newunicodechar{≪}{\ensuremath{\ll}}
\newunicodechar{↛}{\ensuremath{\nrightarrow}}
\newunicodechar{・}{\textperiodcentered}
\newunicodechar{⇒}{\ensuremath{\Rightarrow}}
\newunicodechar{∪}{\ensuremath{\cup}}
\newunicodechar{≥}{\ensuremath{\geq}}
\newunicodechar{≤}{\ensuremath{\leq}}
\newunicodechar{⟨}{\ensuremath{\langle}}
\newunicodechar{⟩}{\ensuremath{\rangle}}
\usepackage{bookmark}
\IfFileExists{xurl.sty}{\usepackage{xurl}}{} % add URL line breaks if available
\urlstyle{same}
% fallback for those not using the hyperref driver hyperxmp:
\makeatletter
\@ifundefined{xmpquote}{\newcommand{\xmpquote}[1]{#1}}{}
\makeatother
\hypersetup{
  pdftitle={The Reality Code Engine: One Form{,} Many Fits{,} and the Low-Entropy Path to Autonomous Coding},
  pdfauthor={David Charlot, Open Interface Engineering},
  colorlinks=true,
  linkcolor={blue},
  filecolor={Maroon},
  citecolor={Blue},
  urlcolor={blue},
  pdfcreator={LaTeX via pandoc}}

\title{The Reality Code Engine: One Form, Many Fits, and the Low-Entropy
Path to Autonomous Coding}
\author{David Charlot, Open Interface Engineering}
\date{Version 1.1, 11 October 2026. HTML:
https://research.openie.dev/papers/rce/}

\begin{document}
\maketitle

{
\setcounter{tocdepth}{2}
\tableofcontents
}
\emph{Version 1.1, 11 October 2026. Data and code:
\href{https://research.openie.dev/data/\#rce}{research.openie.dev/data/}.
LaTeX PDF:
\href{https://research.openie.dev/pdfs/latex/rce.pdf}{/pdfs/latex/rce.pdf}
(\href{https://research.openie.dev/pdfs/latex/rce.tex}{source}).}

\emph{One Form. Many fits. The gate is correctness, and the ranking is
joules.}

The economic reading of this paper is in its companion,
\href{https://research.openie.dev/papers/ai-coding-market/}{The Promised
Trillion Dollar AI-Coding Market}. This paper holds the engineering and
the theory only.

\subsection{1. Thesis}\label{thesis}

There are two paths to autonomous computer automation. They are told
apart by two quantities: the energy spent per correct result, and the
entropy the system must resolve at run time.

The first path generates. A large language model (LLM), a statistical
text model, writes each change from the prompt. Every task pays for
generation again, and every output is a fresh sample from a high-entropy
distribution.

The second path converges. People have already coded most of what will
ever be coded. A program is a database functional state machine: one
Form, many fits. Human intent is high entropy. The machine floor is low
entropy. Code converges by design, so most of the work is to retrieve a
fit that exists, adapt it, and prove it. The Reality Code Engine (RCE)
is built on the second path.

Both paths face the same gate. The gate is correctness: a task's
completeness predicate \(C(z)\), here the task's tests, a proof or an
agreement trace. Among runs that pass the gate, systems are ranked by
joules to completion, \(J^*\). Correctness is a constraint, not a score.

Mixture of Limits is the philosophy: there is a minimal amount of
anything to accomplish a task, and any more is waste
(\href{https://research.openie.dev/papers/gates/\#part-ii-the-cascade-mixture-of-limits}{Metered
Commit Gates, Part II}). RCE is a Mixture of Limits instance. Its gears
run cheapest first, and the first gear that passes the gate closes the
task.

flowG is a stand-alone language project. RCE is its academic first-pass
demonstration. The two stay separate projects.

Energy figures are labeled by source. \texttt{est\_j} is an analytical
estimate. \texttt{reported\_j} is an operating-system or on-chip
counter. \texttt{measured\_j} is reserved for a wall meter and is empty
in every result here.

\subsection{2. Two paths to autonomy}\label{two-paths-to-autonomy}

\begin{figure}
\centering
\includegraphics[width=0.95\linewidth,height=\textheight,keepaspectratio,alt={One
sentence in. One floor out. Human intent is high entropy; a funnel of
conversions from intent to spec to 626 surface languages to an
intermediate representation to machine code loses information at each
boundary; all code reaches a few machine shapes; RCE writes the meaning
once in one canonical Form, treats every surface as a projection, and
lets the cheapest correct form win.}]{fig/intent-entropy.png}
\caption{One sentence in. One floor out. Human intent is high entropy; a
funnel of conversions from intent to spec to 626 surface languages to an
intermediate representation to machine code loses information at each
boundary; all code reaches a few machine shapes; RCE writes the meaning
once in one canonical Form, treats every surface as a projection, and
lets the cheapest correct form win.}
\end{figure}

\emph{Figure 1. Intent entropy and the machine floor. A
\href{https://research.openie.dev/figures/rce/intent-entropy.png}{PNG
copy} is provided. Sources: Shannon (1951); Lopes et al.~(2017); the
626-language count is David Charlot's lowering-path survey, which is
unpublished; its list and method are to be published with this paper's
data, and until then the count is the author's. The Python and C results
are language specification behavior; C signed overflow is undefined
behavior.}

\textbf{Intent is high entropy.} ``Make the checkout faster'' carries
many meanings: shorter page load, fewer steps, faster card approval, a
cached cart, or faster on one device class. Context is implied and never
written. Shannon measured printed English at about 0.6 to 1.3 bits per
letter {[}1{]}. The letters are predictable. The meaning is not.

\textbf{Every conversion loses information.} Intent becomes a
specification, and context is lost. The specification becomes one of
hundreds of surface languages, and alternatives and edge cases are lost.
The source becomes an intermediate representation (IR), and names, style
and the author's reasons are lost. The IR becomes machine code, and
types become sizes and loops become jumps.

\textbf{The floor is low entropy.} Every program ends as loads, stores,
arithmetic and jumps. The destination has room for only a few shapes, so
code bends toward them. Convergence is by design.

\textbf{Boundaries are where meaning diverges.} The same line gives
different answers at the edges. In Python, \texttt{-3\ \%\ 2} is 1 and
\texttt{-7\ /\ 2} is -3.5. In C, \texttt{-3\ \%\ 2} is -1,
\texttt{-7\ /\ 2} is -3, and signed overflow is undefined. A system that
regenerates code at every boundary guesses at every boundary.

\textbf{The two paths, stated as quantities.}

{\def\LTcaptype{none} % do not increment counter
\begin{longtable}[]{@{}
  >{\raggedright\arraybackslash}p{(\linewidth - 6\tabcolsep) * \real{0.2500}}
  >{\raggedright\arraybackslash}p{(\linewidth - 6\tabcolsep) * \real{0.2500}}
  >{\raggedright\arraybackslash}p{(\linewidth - 6\tabcolsep) * \real{0.2500}}
  >{\raggedright\arraybackslash}p{(\linewidth - 6\tabcolsep) * \real{0.2500}}@{}}
\toprule\noalign{}
\begin{minipage}[b]{\linewidth}\raggedright
Path
\end{minipage} & \begin{minipage}[b]{\linewidth}\raggedright
What runs per task
\end{minipage} & \begin{minipage}[b]{\linewidth}\raggedright
Entropy resolved at run time
\end{minipage} & \begin{minipage}[b]{\linewidth}\raggedright
Energy per correct result
\end{minipage} \\
\midrule\noalign{}
\endhead
\bottomrule\noalign{}
\endlastfoot
Generate & A model samples a new change & The full distribution of the
prompt & Paid again on every task \\
Converge (RCE) & Lookup, rule, solver; a model only last & Only the
residual that no stored fit covers & Paid once per new case; then the
handler floor \\
\end{longtable}
}

RCE writes the meaning once, in one canonical Form. Every surface
language is a projection of that Form, checked to match it. Among
equivalent forms, the cheapest correct one runs.

\subsubsection{2.1 Programming is coding
theory}\label{programming-is-coding-theory}

The thesis of this section is David Charlot's: all of programming is
coding theory. A programmer converts intent across two entropic regimes.
Intent is high entropy, a distribution over meanings. Code is low
entropy, one deterministic transition per state. Every step of software
work, from a requirement to a register, is a step of that conversion.
The literature has said so in pieces for seventy years, in many
languages. This section assembles the pieces into one argument and
places each RCE gear on it.

\textbf{The measure.} Shannon defined the entropy of a source and the
capacity of a channel {[}40{]}, and measured printed English at about
0.6 to 1.3 bits per letter {[}1{]}. Kolmogorov, writing in Russian in
\emph{Problemy Peredachi Informatsii} in 1965, gave three approaches to
the quantity of information and made the algorithmic one primary: the
information in an object is the length of the shortest program that
produces it {[}41{]}. Solomonoff had reached the same object as a theory
of inductive inference {[}42{]}, and Chaitin showed that program size
obeys identities formally identical to Shannon's {[}43{]}. Rissanen
turned the idea into a rule for choosing models, the minimum description
length {[}44{]}. Together they say that a program is a code word: a
description of a behavior, and the shortest correct one is the floor.
RCE's extraction rule, keep the least-joule fit among equivalent fits
(Proposition 3), is that rule with joules in place of bits.

\textbf{Code sits in the low-entropy regime, and people measured it.}
Hindle, Barr, Su, Gabel and Devanbu found that code is ``very
repetitive, and in fact even more so than natural languages,'' at lower
cross-entropy under the same n-gram models {[}35{]}. Gabel and Su found
that most lines of a project already exist elsewhere {[}34{]}, and Lopes
et al.~found that most files on GitHub are clones {[}33{]}. Tu, Su and
Devanbu found that code is also local: what was written nearby predicts
what comes next {[}46{]}. Ray and colleagues found that buggy lines are
less natural, with higher entropy, and that the entropy falls when the
bug is fixed {[}45{]}. In 2025 Yang, Kolak, Hellendoorn and Martins used
the entropy change a patch induces to rank repair candidates and to
separate correct patches from overfitting ones {[}51{]}. Allamanis,
Barr, Devanbu and Sutton surveyed the field under the name naturalness
{[}47{]}. The result is not only an English-language one. Prem Devanbu
stated it as a keynote at the India Software Engineering Conference in
2013 {[}50{]}. Sridhara, Sinha and Mani at IBM Research India measured
the natural-language artifacts around code at the same venue and found
that commit messages are the most regular and Stack Overflow prose
behaves like ordinary English {[}49{]}. In China, the \emph{Journal of
Software} (软件学报) published a review of code naturalness research in
2022, which states the cross-entropy measure and the finding that
defective code has higher entropy {[}48{]}. In Japan, Doi, Higo, Arima,
Shimonaka and Kusumoto at Osaka University used n-gram naturalness,
first in a Japanese-language IEICE technical report and then at ICPC
2018, to tell machine-generated source files from human-written ones
{[}73, 74{]}. Generated code sits even lower on the entropy scale than
human code, because a generator is a fixed conversion. In Brazil,
Arrebola's doctoral thesis built code suggestion on n-gram and topic
models segmented by project and maintainer, and evaluated it by
cross-entropy, the bits needed to encode the next token {[}75{]}. The
two regimes are visible in one repository: the prose around the code is
high entropy, and the code is low.

\textbf{The measure matters.} Febres, Jaffé and Gershenson, working from
Venezuela and Mexico, compared English, Spanish and program text by
vocabulary diversity per message length and found the program text more
diverse {[}76{]}. Rahman, Palani and Rigby found that much of the n-gram
predictability of code comes from syntax tokens, and that code is still
repetitive, less than first reported, once those are removed, with API
usage the most repetitive part {[}77{]}. Both results sharpen the claim
rather than reverse it. Diversity of symbols is not predictability, and
the low-entropy regime is a property of the conditional distribution of
the next step, which is what a deterministic transition fixes. RCE's G0
keys on shapes, not on surface tokens, for that reason.

\textbf{Intent sits in the high-entropy regime.} Zave and Jackson showed
that a requirement is a statement about the environment, a specification
is a statement at the interface, and the gap between them is filled only
by domain knowledge that is often unwritten {[}64{]}. Brooks split the
difficulty of software into essence and accident: the essence is the
conceptual construct, and most of it is specification and design, not
representation {[}61{]}. Naur argued that a program is the visible part
of a theory held by its programmers, and that the theory is not fully
recoverable from the text {[}62{]}. Read through Shannon, all three say
the same thing. Intent carries more uncertainty than any one program,
and writing the program is the act that removes it. What remains in a
programmer's head after the code is written is the part of the intent
the code did not encode.

\textbf{The conversion is a refinement, and refinement is a narrowing.}
Hoare gave programs a logic in which a program is correct relative to a
pre- and postcondition {[}54{]}. Wirth taught stepwise refinement:
decompose a task into instructions by successive decisions, each one
fixing a choice the specification left open {[}53{]}. Dijkstra's guarded
commands made the open choices explicit as nondeterminacy, and derived
programs from specifications by removing it {[}52{]}. In Munich, Bauer
and Wössner's CIP project built program development as a sequence of
correctness-preserving transformations from specification to program
{[}58{]}. In France, Abrial's B method and then Event-B carried
refinement to industrial scale, with proof obligations at each step
{[}55, 56{]}. In Finland, Back and von Wright gave refinement its
lattice-theoretic calculus {[}57{]}. In each, a specification denotes a
set of acceptable behaviors and each refinement step shrinks the set,
until one deterministic behavior remains. That is entropy reduction with
a proof at every step. RCE's Form is the last element of that chain: a
total, deterministic transition \(\delta\) over database states, with
every choice resolved
(\href{https://research.openie.dev/papers/rce/\#32-definitions}{§3.2}).

\textbf{Synthesis is search that narrows a hypothesis space.} Gulwani's
FlashFill learns a string program from a few input-output examples by
representing the set of all consistent programs and intersecting it with
each new example {[}59{]}. Gulwani, Polozov and Singh frame program
synthesis in general as search through a space of programs constrained
by intent {[}60{]}, and Solar-Lezama's sketching fills holes in a
partial program with a solver {[}27{]}. Each example or constraint is a
measurement that removes hypotheses. Shannon's measure counts how many
such measurements are needed. RCE's G2 is that gear.

\textbf{What the conversion produces has one shape.} Gurevich proved
that every sequential algorithm is captured by an Abstract State Machine
at its own level of abstraction {[}2{]}. Codd made the state a relation
{[}5{]}, and event sourcing makes the state the fold of an append-only
log {[}14{]}. Under Curry and Howard, a proof is a program and a
proposition is a type {[}18, 72{]}, so a verified fit is a proof of its
specification. These are the reasons the low-entropy regime has few
shapes: the destination is a state machine over a database, whatever the
surface language.

\textbf{Low entropy decays without a gate.} Lehman observed over decades
of large systems that a program in use must keep changing, and that its
complexity grows unless work is done to reduce it {[}63{]}. Torres,
Baltes, Treude and Wagner measured structural and textual entropy over
the commit histories of 25 projects and found the two rise together
{[}78{]}. Entropy flows back into code from the environment. A system
that keeps code in the low-entropy regime needs a gate on every change,
and RCE's gate is correctness
(\href{https://research.openie.dev/papers/rce/\#43-the-gate}{§4.3}).

\textbf{Compilers already do source coding.} Franz and Kistler shipped
programs as compressed abstract syntax trees and generated native code
on load, which is source coding of a program over its own grammar
{[}67{]}. Fraser inferred statistical models of code from a corpus to
compress it {[}68{]}. Futamura showed in 1971, in the Japanese journal
\emph{Systems, Computers, Controls}, that specializing an interpreter to
a program yields a compiled program, and that specializing the
specializer yields a compiler {[}70{]}. Ershov in Novosibirsk developed
the same idea as mixed computation {[}71{]}. The Cousots' abstract
interpretation computes a sound, lower-information description of every
behavior of a program {[}69{]}. The analogy to channel coding is this
paper's reading, not a cited theorem: types, tests and proofs add
checked redundancy so that a wrong program is detected before it runs,
as a parity check detects a corrupted word.

\textbf{Erasure is where the joules go.} Landauer showed that a
logically irreversible step, the erasure of one bit, dissipates at least
\(k_B T \ln 2\) {[}65{]}. Bennett showed that computation can in
principle be made logically reversible, so the erasure, not the
computation, carries the bound {[}66{]}. Converting high-entropy intent
to one low-entropy program discards every alternative the intent
allowed. That discard is an erasure of information about the
alternatives, and it has a thermodynamic floor. The
\href{https://research.openie.dev/papers/agency/\#1-thesis}{agency
paper} prices each confirmed bit against that floor. RCE pays the
discard once per new case and stores the result, so a repeat case erases
nothing new.

\subsubsection{2.2 How RCE converts, and where each gear
sits}\label{how-rce-converts-and-where-each-gear-sits}

Read as coding theory, a task arrives as a message whose remaining
uncertainty is what no stored fit already resolves. Each gear removes
uncertainty at a different price, and the gears are ordered by that
price
(\href{https://research.openie.dev/papers/rce/\#42-the-gears}{§4.2}).

{\def\LTcaptype{none} % do not increment counter
\begin{longtable}[]{@{}
  >{\raggedright\arraybackslash}p{(\linewidth - 4\tabcolsep) * \real{0.3333}}
  >{\raggedright\arraybackslash}p{(\linewidth - 4\tabcolsep) * \real{0.3333}}
  >{\raggedright\arraybackslash}p{(\linewidth - 4\tabcolsep) * \real{0.3333}}@{}}
\toprule\noalign{}
\begin{minipage}[b]{\linewidth}\raggedright
Gear
\end{minipage} & \begin{minipage}[b]{\linewidth}\raggedright
What it does to the uncertainty
\end{minipage} & \begin{minipage}[b]{\linewidth}\raggedright
Lineage in the argument above
\end{minipage} \\
\midrule\noalign{}
\endhead
\bottomrule\noalign{}
\endlastfoot
G0 Lookup & Recognizes that the intent was already encoded; resolves it
by retrieval, at the handler floor & Repetitiveness and clones {[}33,
34, 35{]}; localness {[}46{]} \\
G1 Rules & Applies a known guarded rewrite; resolves a recurring
residual with no search & Correctness-preserving transformation {[}52,
58{]}; refinement with proof obligations {[}55, 56{]} \\
G2 Synthesis & Searches a hypothesis space and narrows it by examples,
constraints and a solver & Program synthesis as search {[}27, 59,
60{]} \\
G3 to G5 Models & Samples from a learned distribution over code when no
lower gear closes the gap; a proposer only & Cross-entropy models of
code {[}35, 47{]} \\
Gate & Accepts a candidate only with tests, coverage and a proof
certificate & Hoare logic {[}54{]}; Curry-Howard {[}18, 72{]} \\
Promotion & Stores the accepted fit, so the same residual costs G0 next
time & Minimum description length {[}44{]}; Landauer erasure paid once
{[}65{]} \\
\end{longtable}
}

The two paths of the table in Section 2 differ in where the conversion
is paid. The generation path samples the full distribution of the prompt
on every task, so it pays the high-entropy side again each time. The
convergence path pays only for the residual, and only once. The economic
reading of that difference is in
\href{https://research.openie.dev/papers/ai-coding-market/\#4-the-theory-that-prices-the-task}{The
Promised Trillion Dollar AI-Coding Market, Section 4}. The stop rule
behind the gear order is
\href{https://research.openie.dev/papers/gates/\#part-ii-the-cascade-mixture-of-limits}{Mixture
of Limits}, and the per-task stop is
\href{https://research.openie.dev/papers/satiation/}{Task Satiation}.

\subsubsection{2.3 Per-rung ladder
results}\label{per-rung-ladder-results}

A coding task is measured here by the \textbf{intent bits} it needs: the
bits that must be supplied from outside what the engine already holds
(the program around the edit, the libraries, the tests, the index of
known fixes). Writing \(H_{task}\) for the bits that pin the target
artifact down and \(H_{corpus}\) for the bits the corpus already fixes,
the supplied bits are \(I = H_{task} - H_{corpus}\). A gear that ranks
the right edit at rank \(r\) among its candidates has supplied
\(\log_2 r\) bits of search, and the gate supplies the rest. The ladder
orders tasks by edit size, an objective proxy for \(I\): one token (L0),
one statement (L1), one hunk (L2), several hunks in one file (L3),
several files (L4), a function from a specification or tests (L5), an
issue-level change (L6) and a program from a specification (L7). The
rungs, corpora, splits and hypotheses were preregistered before the full
runs; the files are in the
\href{https://github.com/openIE-dev/openie-research-experiments/tree/main/rce/bench/ladder/}{public
repository}.

{\def\LTcaptype{none} % do not increment counter
\begin{longtable}[]{@{}
  >{\raggedright\arraybackslash}p{(\linewidth - 14\tabcolsep) * \real{0.1250}}
  >{\raggedright\arraybackslash}p{(\linewidth - 14\tabcolsep) * \real{0.1250}}
  >{\raggedright\arraybackslash}p{(\linewidth - 14\tabcolsep) * \real{0.1250}}
  >{\raggedright\arraybackslash}p{(\linewidth - 14\tabcolsep) * \real{0.1250}}
  >{\raggedright\arraybackslash}p{(\linewidth - 14\tabcolsep) * \real{0.1250}}
  >{\raggedright\arraybackslash}p{(\linewidth - 14\tabcolsep) * \real{0.1250}}
  >{\raggedright\arraybackslash}p{(\linewidth - 14\tabcolsep) * \real{0.1250}}
  >{\raggedright\arraybackslash}p{(\linewidth - 14\tabcolsep) * \real{0.1250}}@{}}
\toprule\noalign{}
\begin{minipage}[b]{\linewidth}\raggedright
Rung
\end{minipage} & \begin{minipage}[b]{\linewidth}\raggedright
Held-out test corpus
\end{minipage} & \begin{minipage}[b]{\linewidth}\raggedright
N
\end{minipage} & \begin{minipage}[b]{\linewidth}\raggedright
Solved
\end{minipage} & \begin{minipage}[b]{\linewidth}\raggedright
Rate
\end{minipage} & \begin{minipage}[b]{\linewidth}\raggedright
Wilson 95\% CI
\end{minipage} & \begin{minipage}[b]{\linewidth}\raggedright
Median search bits
\end{minipage} & \begin{minipage}[b]{\linewidth}\raggedright
Model calls
\end{minipage} \\
\midrule\noalign{}
\endhead
\bottomrule\noalign{}
\endlastfoot
L0 one token & PIT-style mutants {[}109{]} of 9 real repositories at
pinned commits (Java: commons-math, commons-collections, jsoup,
joda-time, gson; Python: click, attrs, pydantic, sympy) & 520,821 &
520,541 & 0.9995 & 0.9994 to 0.9995 & 2.3 & 0 \\
L0 one token, repositories never tuned on & commons-text, commons-io,
more-itertools, httpx & 23,015 & 23,015 & 1.0000 & 0.9998 to 1.0000 &
2.0 & 0 \\
L1 one statement, Java & ManySStuBs4J 2.0 single-hunk fixes {[}107{]} &
24,072 & 9,288 & 0.386 & 0.380 to 0.392 & 3.5 & 0 \\
L1 one statement, Python & TSSB-3M v1.0.0, shards 14 and 15 {[}108{]} &
157,089 & 27,266 & 0.174 & 0.172 to 0.175 & 4.8 & 0 \\
\end{longtable}
}

L2 to L4 are measured on the Defects4J held-out set of 288 bugs
(\href{https://research.openie.dev/papers/rce/\#81-defects4j}{Section
8.1}). L5 to L7 follow the protocols of
\href{https://research.openie.dev/papers/rce/\#82-swe-bench-lite-offline}{Section
8.2}. Section 6 gives the method, the gate and the energy of each row.

The rates read as intent bits. At L0 the missing bits are almost all on
the line, and a prior-ordered single-token neighbourhood supplies them
in a median of 2.3 bits. At L1 the engine sees only the 7-line hunk. It
closes the categories whose bits are inside the hunk (Java: changed
operator 91\%, swapped boolean literal 84\%, swapped arguments 82\%;
Python: changed boolean literal 98.8\%, swapped arguments 99.4\%) and
leaves open those whose bits are outside it: method names (22\%) and
identifiers (37\%) need the file and its types, and new conjuncts (3\%
to 7\%) and whole-statement rewrites (Python, 1.2\% of a category that
is 34\% of the corpus) are new intent. New intent is what a supplier (a
specification, a test, a person or a local model) provides once, after
which the fit is in the index.

\subsection{3. The Form}\label{the-form}

\subsubsection{3.1 The strongest formal statement of ``same
form''}\label{the-strongest-formal-statement-of-same-form}

The scope of the theorem is stated first. Gurevich's sequential
postulates cover sequential algorithms; parallel and interactive
algorithms need the further postulates of Blass and Gurevich {[}3{]}.
Blass, Dershowitz and Gurevich argue that no single equivalence relation
on algorithms fits every purpose {[}79{]}, and Moschovakis and Vardi set
out the debate over what an algorithm is {[}80, 81{]}. A web handler
that does input and output sits outside the sequential theorem. The Form
handles interaction the way the interactive line does: each input is an
event in \(E\), and the state machine steps on it. So ``one Form'' is a
modelling choice grounded in the sequential theorem and extended by
convention, with production precedents in the Elm architecture {[}83{]},
durable workflow event histories {[}84{]} and stream semantics such as
Flo {[}82{]}. The propositions of Section 3.4 are stated for that
choice.

Gurevich's Abstract State Machine (ASM) theorem is the strongest formal
statement that every algorithm has one shape {[}2{]}. An abstract state
machine is a machine whose states are first-order structures and whose
step is a finite set of guarded updates. Gurevich states three
postulates. Sequential Time: an algorithm has states, initial states and
a one-step map \(\tau\). Abstract State: states are first-order
structures over one fixed vocabulary, closed under isomorphism. Bounded
Exploration: a finite set of terms decides each step. His Main Theorem
(Theorem 6.13) reads: ``For every sequential algorithm A, there exists
an equivalent sequential abstract state machine B.'' Blass and Gurevich
extend it to parallel algorithms {[}3{]}. Dershowitz and Gurevich add a
fourth axiom, that initial operations are computable, and from it derive
Church's thesis {[}4{]}.

\subsubsection{3.2 Definitions}\label{definitions}

\textbf{Definition 1 (the Form).} A program is a tuple

\[P = (S, E, \delta, D, \lambda)\]

\begin{itemize}
\tightlist
\item
  \(S\), the state space: structures over a fixed vocabulary \(\Sigma\).
  Each state is a database instance in Codd's sense {[}5{]}.
\item
  \(E\), the event alphabet, including the logical clock tick.
\item
  \(\delta : S \times E \to S\), the transition: total, deterministic
  and pure, written as a finite set of guarded updates ``if \(\varphi\)
  then \(f(\bar t) := t_0\)'', each step reading a bounded set of
  locations.
\item
  \(D = (s_0, L)\), the durable store: an initial state and an
  append-only event log \(L = e_1, \dots, e_n\).
\item
  \(\lambda : S \times E \to O\), the output map. A Mealy machine reads
  both arguments {[}6{]}; a Moore machine reads the state only {[}7{]}.
\end{itemize}

\textbf{Definition 2 (Form and fit).} The \emph{form} is the tuple shape
with four invariants: \(\delta\) is deterministic and pure; exploration
is bounded; the vocabulary is fixed; state is the fold of the log,

\[\mathrm{state}(D) = \mathrm{foldl}(\delta, s_0, L).\]

The \emph{fit} is the data that fills the shape: the vocabulary
\(\Sigma\), the alphabets \(E\) and \(O\), the rule program \(\Pi\) that
defines \(\delta\), the output rules that define \(\lambda\), and the
initial data \(s_0\). In Gurevich's terms the fit is the vocabulary, the
ASM program and the initial structures. In Codd's terms \(\Sigma\) is
the schema and \(s_0\) is the instance. In product-line terms a fit is a
feature selection {[}8{]}. In syntax-guided synthesis (SyGuS) terms a
fit is a term from a grammar that meets a specification {[}9{]}.

\textbf{Definition 3 (fit equivalence).} Two fits are equivalent when,
for every log \(L\), their output sequences under \(\lambda\) are equal.

\textbf{Definition 4 (fit language).} A grammar whose terms denote fits.
Programming languages, SCXML documents, XState configurations and
\texttt{.rce} specifications are fit languages. A fit language is
complete when every ASM program over \(\Sigma\) has a term in it.

\textbf{Definition 5 (control projection).} Split
\(\Sigma = \Sigma_c \cup \Sigma_d\) into finite control symbols and data
relations. The statechart of \(P\) is the projection of \(\delta\) onto
\(\Sigma_c\): states are values of \(\Sigma_c\), transitions carry an
event and a guard over \(\Sigma_d\), and actions are updates to
\(\Sigma_d\). Harel's statecharts give the hierarchy and the orthogonal
regions {[}10{]}.

\textbf{Definition 6 (fit e-graph).} For a fit \(F\) and a set \(R\) of
sound rewrites, the fit e-graph \(EG(F, R)\) is the saturated e-graph of
\(F\)'s rule terms under \(R\) {[}11, 12{]}. Each e-class holds fits
equivalent by Definition 3.

\textbf{Definition 7 (joule cost).} \(\mathrm{cost}_J(t)\) is the
estimated energy (\texttt{est\_j}) to run term \(t\) once on a named
substrate, weighted by the expected number of runs.

\subsubsection{3.3 Why ``database functional state
machine''}\label{why-database-functional-state-machine}

A finite relational structure is a database instance, so a Gurevich
state is a database and each step is a bounded transaction.
\emph{Database}: \(S\) is a relational structure. \emph{Functional}:
\(\delta\) and \(\lambda\) are pure. \emph{State machine}: runs are
folds of \(\delta\) over events. Moseley and Marks draw the same line in
architecture {[}13{]}. Event sourcing draws it in application state:
``Capture all changes to an application state as a sequence of events''
{[}14{]}. Dedalus draws it in logic, where mutable state is a relation
indexed by time {[}15{]}.

\subsubsection{3.4 Propositions}\label{propositions}

\textbf{Proposition 1 (the Form captures every sequential algorithm).}
For every sequential algorithm \(A\) there is a Form instance \(P\)
whose runs are the runs of \(A\).

\emph{Proof sketch.} Theorem 6.13 gives an ASM \(B\) with the states,
initial states and one-step map of \(A\). Set \(E = \{\varepsilon\}\),
\(\delta(s, \varepsilon) = \tau_B(s)\), \(D = (s_0, \varepsilon^n)\) and
\(\lambda\) the identity. The fold after \(n\) ticks is
\(\tau_A^n(s_0)\). By Gurevich's Corollary 3.3, \(P\) and \(A\) have the
same runs. ∎

\textbf{Proposition 1b (interactive programs).} Extend \(\Sigma\) with a
nullary input symbol set to \(e\) before each step; then
\(\delta(s, e) = \tau(s[\mathit{in} := e])\) is a sequential algorithm
over the extended vocabulary, and Proposition 1 applies. The
construction is ours.

\textbf{Proposition 1c (computability).} Under the Dershowitz and
Gurevich axiom, every Form instance computes a Turing-computable
function {[}4, 16, 17{]}. The Form adds no power. It gives every program
in the class one shape.

\textbf{Proposition 1d (proofs are fits).} ``For each proof of a given
proposition, there is a program of the corresponding type'' {[}18{]}. A
type is a specification and a proof is a fit that meets it. The
correctness gate is the existence of that proof or of a passing
acceptance test.

\textbf{Proposition 2 (one generator).} One function
\(G : \mathrm{Fit} \to \mathrm{Program}\) generates a runnable program
for every fit and reads no domain. \emph{Basis:} every program is a Form
instance, so \(G\) emits a typed store, an event type, compiled rules, a
log with a fold and an output dispatcher. Scaffolders such as Rails
generators, JHipster and Yeoman are \(G\) restricted to small fit
languages. Product-line engineering starts from ``the systematic
discovery and exploitation of commonality across related software
systems'' {[}8{]}; the commonality is the Form.

\textbf{Proposition 3 (Mixture of Limits as extraction).} Extraction
with \(\mathrm{cost}_J\) returns the least-joule term the extractor
finds in the e-class of \(F\). \emph{Proof sketch.} The rewrites are
sound, so every term in the e-class has \(F\)'s observable behavior and
passes the same gate. Extraction minimizes the cost, and the cost is
joules. ∎ This is Mixture of Limits applied to code: keep the minimal
fit that does the task, and the rest is waste.

\textbf{Proposition 4 (replay determinism).} If \(\delta\) reads no wall
clock, orders updates without names and logs every input, two replays of
one \(D\) give identical state hashes at every step.

\textbf{Proposition 5 (statechart projection is sound).} Project each
state of a run onto \(\Sigma_c\). The result is a run of the emitted
statechart under the same events.

\textbf{Proposition 6 (coordination budget).} Rules tagged monotone run
on replicas without coordination; only non-monotone rules need a
coordination point. \emph{Basis:} the CALM theorem, ``the programs that
have consistent, coordination-free distributed implementations are
exactly the programs that can be expressed in monotonic logic''
{[}19{]}.

\textbf{Proposition 7 (intelligence as a joule derivative).} Let \(X\)
count tasks through the gate and \(J\) the joules spent.
\(\iota = dX/dJ\) is read in hindsight from \(D\), which records every
event and its \texttt{est\_j}. Extraction under \(\mathrm{cost}_J\)
holds the gate fixed and does not raise \(J\), so at equal \(X\) it does
not lower \(\iota\)
(\href{https://research.openie.dev/papers/agency/\#23-layer-3-calculus-intelligence-as-a-derivative-read-in-hindsight}{Agency,
Section 2.3}).

\subsubsection{3.5 Graph, e-graph and
statechart}\label{graph-e-graph-and-statechart}

{\def\LTcaptype{none} % do not increment counter
\begin{longtable}[]{@{}
  >{\raggedright\arraybackslash}p{(\linewidth - 4\tabcolsep) * \real{0.3333}}
  >{\raggedright\arraybackslash}p{(\linewidth - 4\tabcolsep) * \real{0.3333}}
  >{\raggedright\arraybackslash}p{(\linewidth - 4\tabcolsep) * \real{0.3333}}@{}}
\toprule\noalign{}
\begin{minipage}[b]{\linewidth}\raggedright
Structure
\end{minipage} & \begin{minipage}[b]{\linewidth}\raggedright
Part of the Form
\end{minipage} & \begin{minipage}[b]{\linewidth}\raggedright
What it decides
\end{minipage} \\
\midrule\noalign{}
\endhead
\bottomrule\noalign{}
\endlastfoot
Statechart & Guards of \(\delta\) over \(\Sigma_c\) & Which transition
fires \\
Dataflow graph & Read and write dependencies of \(\delta\) and
\(\lambda\) (Kahn process networks {[}20{]}) & What runs, and in what
order \\
E-graph & The equivalence class of the fit & Which equivalent fit
runs \\
\end{longtable}
}

The statechart and the graph describe one fit. The e-graph describes
every fit equivalent to it. Extraction chooses among them by joules.

\subsection{4. The engine: gears, gate and
receipts}\label{the-engine-gears-gate-and-receipts}

\subsubsection{4.1 The engine is a Form}\label{the-engine-is-a-form}

RCE runs software as a graph of small deterministic handlers. A handler
that cannot decide raises an escalation. A resolver answers the
escalation once, and the answer is installed as a permanent handler, so
the next case of the same shape runs with no resolver. The engine meets
every Form invariant, each checked by a named test in
\texttt{tests/form\_laws.rs}:

{\def\LTcaptype{none} % do not increment counter
\begin{longtable}[]{@{}
  >{\raggedright\arraybackslash}p{(\linewidth - 2\tabcolsep) * \real{0.5000}}
  >{\raggedright\arraybackslash}p{(\linewidth - 2\tabcolsep) * \real{0.5000}}@{}}
\toprule\noalign{}
\begin{minipage}[b]{\linewidth}\raggedright
Invariant
\end{minipage} & \begin{minipage}[b]{\linewidth}\raggedright
What the engine does
\end{minipage} \\
\midrule\noalign{}
\endhead
\bottomrule\noalign{}
\endlastfoot
Reload is a verified fold & Reload reads the log verbatim, recomputes
every hash and link, and folds the rest with \(\delta\). Nothing is
re-stamped. A broken chain refuses to load. \\
Name-free tick order & Handlers run in topological rank order, then by a
name-free node key. Outputs are staged and delivered after all handlers
run, as in ASM semantics. Renaming every node gives the same anonymous
state hash at every entry. \\
No silent drop & A full mailbox holds the message in an ordered queue
and logs \texttt{Backpressure}. With capacity 1 and 10 tasks, 9
\texttt{Backpressure} entries are logged and all 10 tasks close. \\
Payloads are in the log & Inputs, ticks, external handler results,
installs and faults are log entries. The state hash after every entry
matches a fold from an empty engine, on 1,000 random logs. \\
Logical clock & Liveness counts ticks, not wall time. \\
Stable hashes & SHA-256 (Secure Hash Algorithm, 256-bit) over canonical
JSON, keys sorted at every depth, for entries, states, shapes and rule
ids. \\
\end{longtable}
}

\subsubsection{4.2 The gears}\label{the-gears}

Each gear is tried only when the gear below it does not pass the gate.
That order is one fixed cascade, and the fixed cascade is one Mixture of
Limits instance.

{\def\LTcaptype{none} % do not increment counter
\begin{longtable}[]{@{}
  >{\raggedright\arraybackslash}p{(\linewidth - 4\tabcolsep) * \real{0.3333}}
  >{\raggedright\arraybackslash}p{(\linewidth - 4\tabcolsep) * \real{0.3333}}
  >{\raggedright\arraybackslash}p{(\linewidth - 4\tabcolsep) * \real{0.3333}}@{}}
\toprule\noalign{}
\begin{minipage}[b]{\linewidth}\raggedright
Gear
\end{minipage} & \begin{minipage}[b]{\linewidth}\raggedright
Method
\end{minipage} & \begin{minipage}[b]{\linewidth}\raggedright
Lineage
\end{minipage} \\
\midrule\noalign{}
\endhead
\bottomrule\noalign{}
\endlastfoot
G0 Lookup & A shared handler library keyed by SHA-256 shape hash; for
code repair, localization to ranked functions & Code search and
retrieval {[}21{]} \\
G1 Rules & Deterministic, guarded repair templates and schema-derived
validators, each checked for mutation sensitivity & PAR, TBar, AVATAR,
FixMiner {[}22, 23, 24, 25{]}; codemods such as Coccinelle {[}26{]} \\
G2 Synthesis & Enumeration over the handler grammar, solvers and
e-graphs & Sketching, SyGuS, egg {[}27, 9, 12{]} \\
G3 to G5 Models & A typed decision model, a local small model, then a
remote model, last & Model cascades {[}28{]} \\
\end{longtable}
}

\subsubsection{4.3 The gate}\label{the-gate}

A candidate is promoted only when it passes the gate. The gate replays
the node's recorded inputs, runs the candidate in shadow on fresh
samples with a coverage bound (the rule of three at a declared error
budget), checks the output schema, and requires the proof certificate of
Section 5. A rejected candidate is logged with its counterexample.

\subsubsection{4.4 Receipts}\label{receipts}

Every promotion, refusal, deferral and install is a mutation in an
append-only log chained with SHA-256, so tampering with any past entry
breaks the chain {[}29, 30{]}. Each handler's cost is priced from
interpreter operation counts with the host cost table of FGIR, the
native-code backend of flowG, and labeled \texttt{est\_j}. Each foreman
call carries its receipt. The ladder runs add a package energy counter
on a subsample, labeled \texttt{reported\_j}.

\subsubsection{4.5 The Form compiler}\label{the-form-compiler}

\texttt{rce\ form} takes a specification and its fit and emits five
targets from one Form IR:

\begin{verbatim}
.rce spec + fit
  lower ─► normalize ─► egg e-graph, equality saturation (+ admitted proposals)
  ─► extract (cost = est_j) ─► Form IR (typed relations, regions, rules with SHA-256 ids)
     ├─ Rust crate (pure step, typed state, SHA-256 log and fold, tests)
     │  └─ WASM (the same crate on wasm32-wasip1)
     ├─ Python (standard library only)
     ├─ XState v5 machine
     └─ SCXML document, executed by SCION
\end{verbatim}

The IR fixes one behavior at every numeric edge: checked 64-bit integer
overflow is an error and never wraps; integer \texttt{/} truncates
toward zero (\(-7/2 = -3\)); \texttt{\%} takes the dividend's sign
(\(-3 \bmod 2 = -1\)); division by zero is an error; floats are IEEE 754
binary64 and a non-finite result is an error; null and absent are
distinct; time is the logical clock only. Every target is generated.
None is authored.

A process, HTTP or closure handler has no term, so its logged results
compile to a recorded oracle keyed by the SHA-256 of the canonical
input. \texttt{rce\ scan\ DIR\ -\/-form} lifts existing Rust, Python,
TypeScript and Go into the Form's schema, and every lifted element is
labeled as a syntactic lift.

\subsection{5. Verification}\label{verification}

\subsubsection{5.1 Differential testing across five
targets}\label{differential-testing-across-five-targets}

\texttt{tests/form\_targets.rs} drives an eleven-operator arithmetic
specification with 29 edge events plus 500 seeded random edge values.
The generated Rust, Python, WASM, XState and SCXML targets print the
reference trace byte for byte, including the arithmetic edges. This is
differential testing in the line of McKeeman and of Csmith {[}86{]}:
agreement among independent implementations on a stated input set. It is
not translation validation {[}85{]} and not a verified compiler. It
covers the eleven operators and the 529 events in the test, and it says
nothing about inputs outside them. The Z3 and Kani results below are the
verification.

\subsubsection{5.2 Rewrites proved in Z3}\label{rewrites-proved-in-z3}

Verified rewrite rules have a lineage. Ruler infers and checks rewrite
rules by equality saturation with no model in the loop {[}87{]}. The 3
of 16 model-proposed rules admitted in Section 5.5 are therefore a
statement about the gate, which refused 13, and not a claim that a model
is the best rule source.

Every built-in e-graph rewrite is proved in the Z3 SMT solver
(satisfiability modulo theories) {[}31{]} against a quantifier-free
encoding of the Form's term evaluation and arithmetic: values as a
datatype, 64-bit bit-vectors with checked overflow, signed division and
remainder with the dividend's sign, and binary64 floats. A rewrite that
changes overflow, escalation order or a message gets no proof. A
proposed equivalence enters the e-graph only with a proof.

\subsubsection{5.3 Kani proofs of the
kernel}\label{kani-proofs-of-the-kernel}

The engine's kernel (hash chain, fold, name-free positions and handler
evaluation) is checked with the Kani model checker for Rust {[}32{]}.
Handler evaluation is the fixed point of a non-recursive step, and Kani
proves that step total and deterministic per constructor, by structural
induction. Ten harnesses are verified, with 0 of 505 checks failed:

{\def\LTcaptype{none} % do not increment counter
\begin{longtable}[]{@{}
  >{\raggedright\arraybackslash}p{(\linewidth - 2\tabcolsep) * \real{0.5000}}
  >{\raggedright\arraybackslash}p{(\linewidth - 2\tabcolsep) * \real{0.5000}}@{}}
\toprule\noalign{}
\begin{minipage}[b]{\linewidth}\raggedright
Harness
\end{minipage} & \begin{minipage}[b]{\linewidth}\raggedright
Property
\end{minipage} \\
\midrule\noalign{}
\endhead
\bottomrule\noalign{}
\endlastfoot
\texttt{chain\_appended\_verifies} & An appended hash chain verifies \\
\texttt{chain\_detects\_tampering} & A changed entry breaks the chain \\
\texttt{fold\_equals\_replay} & The fold equals replay \\
\texttt{fold\_rejects\_wrong\_state} & A fold that disagrees with the
stored state is refused \\
\texttt{position\_is\_name\_independent} & Rule positions do not depend
on node names \\
\texttt{reason\_tokens\_are\_name\_independent} & Reasons are name-free
by whole tokens \\
\texttt{handler\_core\_leaves} & Leaf handlers are total and
deterministic \\
\texttt{handler\_core\_validate} & Validation is total and
deterministic \\
\texttt{handler\_core\_pipeline} & Pipelines are total and
deterministic \\
\texttt{handler\_core\_ifhas} & Guards are total and deterministic \\
\end{longtable}
}

\subsubsection{5.4 Promotion gate on the proof
certificate}\label{promotion-gate-on-the-proof-certificate}

The Kani run writes a certificate whose hash is stored in
\texttt{proofs/kani.cert}. The engine, the log, rule ids and handler
evaluation route through the proved kernel, and promotion requires the
certificate.

\subsubsection{5.5 The model as a proposer behind the
check}\label{the-model-as-a-proposer-behind-the-check}

A model may propose equivalences. It may not admit them. With a local
1.5-billion-parameter coding model (Qwen2.5-Coder-1.5B under llama.cpp)
on the demonstration specification, the proposer offered 16
equivalences. The check admitted 3, by randomized trials and a Z3 proof,
and rejected 13. A deterministic proposer of 87 familiar identities had
9 admitted, all structural, and every arithmetic identity rejected for
fault order, non-numbers or overflow.

\subsubsection{5.6 Brownfield agreement}\label{brownfield-agreement}

The lift reads the same order machine from Rust, Python, TypeScript and
Go fixtures, runs the original code on 40 seeded sequences of 30 calls,
and agrees on all 1,200 calls per language. A frozen benchmark,
preregistered before any result, re-expressed three existing
sales-outreach programs as Forms. The Forms agree with the original code
on 4,500 of 4,500 recorded inputs, and extraction lowers \texttt{est\_j}
on each (for example \(1.69 \times 10^{-8}\) J to
\(1.49 \times 10^{-8}\) J).

\subsection{6. Results}\label{results}

\textbf{Method.} Each corpus is shuffled with a fixed seed. The first
2,000 positions of each shuffle, and for L0 three whole repositories
(commons-lang, requests, flask), form the development split, on which
the gears may be tuned; development results are not reported. Every
remaining task is in the test split. Before a test run the engine is
committed and its binary hash recorded, and no gear changes during the
run. Every row in this section is a frozen-engine result on test tasks
only, with 0 model calls. The preregistration (\texttt{PREREG.json}),
the per-task rows and the report script are in the
\href{https://github.com/openIE-dev/openie-research-experiments/tree/main/rce/bench/ladder/}{public
repository}.

\textbf{Gate.} At L0 and L1 the location is given and the gate is token
identity: the repaired line or statement must equal the developer's,
token for token. This is stricter than passing the tests, because a
patch can pass a weak test suite and still be wrong. Qi and colleagues
found that most patches GenProg reported as fixes were plausible but
incorrect {[}92{]}; overfitting to tests is documented in search-based
and semantics-based repair {[}93, 94{]} and addressed with generated
tests {[}95{]}; tools score higher on the benchmark they were tuned on
{[}96{]}; and fault localization biases comparisons {[}97{]}. Token
identity removes the plausibility gap on these rungs, and giving the
location removes the localization term, so the rate measures the repair
gears alone. L2 and above use the tests as the gate and report plausible
and correct counts separately (Section 8.1).

\subsubsection{6.1 L0: one-token faults}\label{l0-one-token-faults}

The L0 corpus is every PIT-style mutant {[}109{]} with a valid
replacement token in nine real repositories at pinned commits, five Java
and four Python: 520,821 test tasks out of 547,448. The engine (commit
\texttt{65571d0}, gears G0-T and G1-J) restores 520,541, a rate of
0.9995 (95\% Wilson interval 0.9994 to 0.9995) at a median of 2.3 search
bits. Of the 280 unsolved, 108 are mutated digits inside floating-point
literals and 2 are inside Python set literals. A float-literal gear
(G0-F, commit \texttt{ecfe5a1}) is tested on four repositories never
used for tuning (commons-text, commons-io, more-itertools, httpx):
23,015 of 23,015 tasks, 1.0000 (0.9998 to 1.0000), a median of 2.0
search bits. Mutants are a stand-in for real one-token faults that Just
and colleagues found to be coupled with real faults {[}110{]}.

\subsubsection{6.2 L1: one-statement bugs from real
histories}\label{l1-one-statement-bugs-from-real-histories}

L1 uses two public datasets of real single-statement bug fixes mined
from open-source histories. ManySStuBs4J 2.0 {[}107{]} gives 24,072 Java
test tasks; the engine (commit \texttt{f2701b7}) resolves 9,288, a rate
of 0.386 (0.380 to 0.392). TSSB-3M v1.0.0 {[}108{]} gives 157,089 Python
test tasks from shards 14 and 15; the same engine resolves 27,266, a
rate of 0.174 (0.172 to 0.175). The engine receives the 7-line hunk
only, so the Java rate is a hunk-only result; the full-file result is
the next row of this table. Per-category rates are in Section 2.3 and
the full per-stratum tables are in the repository. These datasets have
no published repair leaderboard under this gate (hunk only, token
identity, location given), so no learned-repair number is set beside
them; Defects4J, which does have one, is compared with TBar in Section
8.1.

\subsubsection{6.3 Energy}\label{energy}

Every row carries \texttt{est\_j}, an estimate of CPU seconds times 10
W: 0.40 J per solved task at L0, 0.13 on the fresh repositories, 0.10 on
ManySStuBs4J and 0.27 on TSSB-3M. A subsample was read on RAPL
{[}105{]}, the processor's on-chip package counter, on a Linux laptop
with an Intel i9-13900H (hostname \texttt{jetson-hub}; not an NVIDIA
Jetson). The harness ran on 4 pinned CPUs, and \texttt{reported\_j} is
the package delta over the whole run, Python driver included, less idle
power from a 10 s sample taken just before, times the run's duration.
Continuous-integration runners were up, so idle was 40 W to 44 W and the
uncertainty is about ±1 W times the duration, ±10 J on these totals.

{\def\LTcaptype{none} % do not increment counter
\begin{longtable}[]{@{}
  >{\raggedright\arraybackslash}p{(\linewidth - 12\tabcolsep) * \real{0.1429}}
  >{\raggedright\arraybackslash}p{(\linewidth - 12\tabcolsep) * \real{0.1429}}
  >{\raggedright\arraybackslash}p{(\linewidth - 12\tabcolsep) * \real{0.1429}}
  >{\raggedright\arraybackslash}p{(\linewidth - 12\tabcolsep) * \real{0.1429}}
  >{\raggedright\arraybackslash}p{(\linewidth - 12\tabcolsep) * \real{0.1429}}
  >{\raggedright\arraybackslash}p{(\linewidth - 12\tabcolsep) * \real{0.1429}}
  >{\raggedright\arraybackslash}p{(\linewidth - 12\tabcolsep) * \real{0.1429}}@{}}
\toprule\noalign{}
\begin{minipage}[b]{\linewidth}\raggedright
Subsample (test split)
\end{minipage} & \begin{minipage}[b]{\linewidth}\raggedright
Tasks
\end{minipage} & \begin{minipage}[b]{\linewidth}\raggedright
Solved
\end{minipage} & \begin{minipage}[b]{\linewidth}\raggedright
Raw J
\end{minipage} & \begin{minipage}[b]{\linewidth}\raggedright
\texttt{reported\_j}
\end{minipage} & \begin{minipage}[b]{\linewidth}\raggedright
\texttt{reported\_j} per solved task
\end{minipage} & \begin{minipage}[b]{\linewidth}\raggedright
Seconds
\end{minipage} \\
\midrule\noalign{}
\endhead
\bottomrule\noalign{}
\endlastfoot
L0, gson mutants (seeded, first 3,000) & 3,000 & 3,000 & 648.3 & 167.5 &
0.056 & 11.8 \\
L1, ManySStuBs4J positions 2,000 to 3,999 & 2,000 & 690 & 377.2 & 17.9 &
0.026 & 8.2 \\
\end{longtable}
}

The L1 subsample builds its project index from its 2,000 tasks only, so
its solve rate (34.5\%) is below the full run's. \texttt{measured\_j}, a
wall meter on a named rail, is empty in every row; its protocol is
Section 8.3.

\subsection{7. Reuse evidence}\label{reuse-evidence}

The second path rests on an engineering fact: code is reused at every
scale. Each method below measures it differently.

{\def\LTcaptype{none} % do not increment counter
\begin{longtable}[]{@{}
  >{\raggedright\arraybackslash}p{(\linewidth - 4\tabcolsep) * \real{0.3333}}
  >{\raggedright\arraybackslash}p{(\linewidth - 4\tabcolsep) * \real{0.3333}}
  >{\raggedright\arraybackslash}p{(\linewidth - 4\tabcolsep) * \real{0.3333}}@{}}
\toprule\noalign{}
\begin{minipage}[b]{\linewidth}\raggedright
Method
\end{minipage} & \begin{minipage}[b]{\linewidth}\raggedright
Source
\end{minipage} & \begin{minipage}[b]{\linewidth}\raggedright
Finding
\end{minipage} \\
\midrule\noalign{}
\endhead
\bottomrule\noalign{}
\endlastfoot
File-level clone mapping & DéjàVu, Lopes et al.~2017 {[}33{]} & Of 428
million files in 4.5 million non-fork GitHub projects, 85 million are
unique: ``70\% of the code on GitHub consists of clones of previously
created files.'' In JavaScript, 6\% of files are distinct. \\
Line-level redundancy & Gabel and Su 2010 {[}34{]} & Against 420 million
lines, ``between 50\% and 100\% of each project is redundant'' at about
one line. \\
Statistical language models & Hindle et al.~2012 {[}35{]} & Code is
``very repetitive, and in fact even more so than natural languages.'' \\
Corpus near-deduplication & The Stack {[}36{]} & Near-deduplication
removes 54.7\% of The Stack v1 (6.4 TB to 2.9 TB) and 52.4\% of v2 (67.5
TB to 32.1 TB). \\
Whole-archive preservation & Software Heritage & About 30 billion source
files from 445 million origins, each stored once by the hash of its
bytes. \\
Mined fix patterns & PAR, TBar {[}22, 23{]} & PAR found ``several common
fix patterns'' in more than 60,000 human patches. TBar's 15 pattern
categories fix 43 Defects4J bugs correctly in the standard pipeline, and
74 of 101 with perfect fault localization, with no learning. \\
Cross-language tasks & Rosetta Code & 1,360 tasks, each one concept
implemented across many of 1,009 known languages. \\
\end{longtable}
}

Each row is an input to the engine. Clones and dedup say a fit usually
exists, which is G0. Recurring fixes say adaptation is mostly a small
catalog of guarded rewrites, which is G1. Line-level redundancy shrinks
the search space of G2. The L0 and L1 rungs measure this at scale on
held-out splits (Section 6). Two cautions sharpen the reading.
Duplicates inflate retrieval and model scores unless removed {[}104{]}.
And buggy lines are less natural than fixed ones {[}45{]}, so the bug
itself is the case lookup is least able to predict; the fix shape, not
the buggy line, is what G0 and G1 must key on. The rule-mining line that
does this at scale, from GenProg's generate-and-validate search {[}88{]}
and Prophet's learned correctness model {[}89{]} to Getafix at Facebook
{[}90{]} and Sorald's model-free repair of static-analysis violations
{[}91{]}, derives templates from disjoint history and reports precision
on unseen cases {[}98{]}. That is the standard Section 8.4 adopts.

\subsection{8. Measurement protocols for what comes
next}\label{measurement-protocols-for-what-comes-next}

Each protocol is fixed before its run: preregistration committed, seeds
and solver commit recorded, every artifact kept. Gate first, joules
second, labels always.

\subsubsection{8.1 Defects4J}\label{defects4j}

Defects4J holds 854 active real Java bugs, each with reproducing tests
{[}37{]}.

\begin{enumerate}
\def\labelenumi{\arabic{enumi}.}
\tightlist
\item
  Develop on a 100-bug subset, offline, through
  \texttt{defects4j\ checkout}; report held-out results only, on a
  disjoint set of 288 bugs frozen before the run
  (\texttt{FROZEN-heldout-v1}).
\item
  Run G0, then G1, then G2, then a local model last. Record the gear
  that closed each bug.
\item
  Gate: the triggering tests and the full suite pass. Report plausible
  and correct counts separately; correct means equivalent to the
  developer patch on review.
\item
  Report model calls per correct fix and joules per correct fix,
  labeled.
\item
  Reference points: TBar's 43 correct fixes and its 74 of 101 under
  perfect localization.
\item
  Contamination check: flag every bug whose fix appears verbatim in G0's
  index, and report with and without flagged bugs.
\end{enumerate}

\subsubsection{8.2 SWE-bench Lite,
offline}\label{swe-bench-lite-offline}

SWE-bench Lite is a 300-task subset of real GitHub issues resolved by
hidden tests {[}38{]}.

\begin{enumerate}
\def\labelenumi{\arabic{enumi}.}
\tightlist
\item
  Use the official Docker harness with no network.
\item
  First run: G0 and G1 only, zero model calls. This is the floor of
  solving without a model.
\item
  Second run: add a local model as the last gear.
\item
  Report the resolve rate, model calls per resolved task and joules per
  resolved task, labeled.
\item
  Contamination check as in 8.1. SWE-Bench+ found that 32.67\% of
  successful patches had the solution in the issue text {[}39{]}.
\end{enumerate}

\subsubsection{8.3 A wall meter}\label{a-wall-meter}

\begin{enumerate}
\def\labelenumi{\arabic{enumi}.}
\tightlist
\item
  Attach a wall meter or a shunt meter on the host's supply, logged at a
  stated rate with its clock synchronized to the run log.
\item
  Idle machine: CI runners and other builds stopped, and a quiet window
  gate checked before each block and recorded.
\item
  Subtract a measured idle baseline, and state the method.
\item
  Fill \texttt{measured\_j} from the meter alone. Report
  \texttt{reported\_j} (RAPL) and \texttt{est\_j} beside it in separate
  columns. Never sum across labels.
\item
  Repeat the RAPL subsamples of Section 6.3 under the meter, then the
  Defects4J held-out set.
\end{enumerate}

\subsubsection{8.4 A held-out run with the gears
frozen}\label{a-held-out-run-with-the-gears-frozen}

This is the method of Section 6, extended to new corpora.

\begin{enumerate}
\def\labelenumi{\arabic{enumi}.}
\tightlist
\item
  Freeze G0, G1 and G2 at a tagged commit, record the binary hash, and
  publish the rule set with its authoring timestamps.
\item
  Draw a disjoint held-out slice: the Rust and Python instances of
  Multi-SWE-bench {[}102{]} and BugsInPy {[}103{]}, de-duplicated
  against G0's index {[}104{]}.
\item
  Preregister the slice, the seeds and the hypotheses before the first
  run.
\item
  Report resolved, plausible and byte-identical counts separately, each
  with a 95\% interval.
\end{enumerate}

\subsubsection{8.5 Baselines in the same
table}\label{baselines-in-the-same-table}

\begin{enumerate}
\def\labelenumi{\arabic{enumi}.}
\tightlist
\item
  Model-free: \texttt{cargo\ fix} and clippy {[}106{]}, and Sorald or a
  TBar port {[}91{]}.
\item
  Cheap model: Agentless with a local model {[}99{]}.
\item
  Remote model: an agent with a frontier model, its energy estimated
  from provider disclosures and labeled \texttt{est\_j}.
\item
  Report cost-of-pass {[}101{]} in joules and in dollars for every row,
  with cost controlled as Kapoor and colleagues require {[}100{]}.
\end{enumerate}

\subsubsection{8.6 Convergence of model calls at large
N}\label{convergence-of-model-calls-at-large-n}

Promotion is the mechanism by which the model leaves the loop: a
resolver is paid once for a new case, the gate admits the fit as a
handler, and every later case of the same shape closes in G0 or G1
(\href{https://research.openie.dev/papers/rce/\#42-the-gears}{Section
4.2}). Convergence is the claim that model calls per round fall to zero
on a fixed specification under live traffic. It is measured as follows.

\begin{enumerate}
\def\labelenumi{\arabic{enumi}.}
\tightlist
\item
  Draw at least 1,000 specifications from public repositories,
  stratified by language and size, and freeze the engine and the
  resolver before the first round.
\item
  Replay recorded production traffic, at least 10,000 events per
  specification, in rounds of fixed size.
\item
  Record per round: new escalations, model calls, promotions, deferrals
  and joules per closed task, each labeled (\texttt{est\_j},
  \texttt{reported\_j}, \texttt{measured\_j}).
\item
  Primary endpoint: the share of specifications reaching a round with
  zero model calls within 20 rounds, with a 95\% Wilson interval.
  Secondary: model calls per closed task by round, and the round of the
  last model call.
\item
  Report specifications that never converge with their cause (a stuck
  input, a missing dependency, new intent), and run stuck-input deferral
  as a preregistered arm.
\end{enumerate}

\textbf{Proposition 8 (floor approach).} As G0 and G1 coverage grows,
joules per solved task approach the gold floor of applying a known fix
and running the tests. The ladder rungs and the Defects4J held-out set
under the wall meter test it.

\subsection{9. Ties to the companion
papers}\label{ties-to-the-companion-papers}

\textbf{Mixture of Limits.}
\href{https://research.openie.dev/papers/gates/\#part-ii-the-cascade-mixture-of-limits}{Metered
Commit Gates, Part II} states the philosophy: a minimal amount of
anything is needed to accomplish a task, and any more is waste. RCE
applies it three times. The gear order runs the cheapest sufficient gear
first. Promotion pays a resolver once per new case, after which every
repeat call would be waste. Extraction keeps the least-joule fit among
equivalent fits (Proposition 3).

\textbf{Metered Commit Gates.}
\href{https://research.openie.dev/papers/gates/}{Part I} says a proposal
may not authorize itself. In RCE a model proposes and the gate commits:
a candidate enters only with tests passed, a coverage bound met and a
proof certificate, and every decision is a receipt in the hash chain.

\textbf{The Universal Law of Agency.}
\href{https://research.openie.dev/papers/agency/}{The agency paper}
defines intelligence as the derivative of agency in joules,
\(\iota = dX/dJ\), read in hindsight. The RCE log is that hindsight
record: \(X\) is tasks through the gate and \(J\) is their joules,
labeled. A deferral with a receipt is the refusal clause of the gate.

\textbf{Task Satiation.}
\href{https://research.openie.dev/papers/satiation/}{Task Satiation
vs.~Aggregate Rebound in AI Demand} says that once a task's completion
predicate is true, more synthesis on that task buys nothing. RCE is
per-task satiation inside one system: once a handler covers a case,
further inference on that case is zero. The aggregate demand question
belongs to that paper and to the
\href{https://research.openie.dev/papers/ai-coding-market/}{economics
companion}.

\subsection{10. Falsifiers}\label{falsifiers}

Fixed before the held-out run of Section 8.4.

{\def\LTcaptype{none} % do not increment counter
\begin{longtable}[]{@{}
  >{\raggedright\arraybackslash}p{(\linewidth - 4\tabcolsep) * \real{0.3333}}
  >{\raggedright\arraybackslash}p{(\linewidth - 4\tabcolsep) * \real{0.3333}}
  >{\raggedright\arraybackslash}p{(\linewidth - 4\tabcolsep) * \real{0.3333}}@{}}
\toprule\noalign{}
\begin{minipage}[b]{\linewidth}\raggedright
ID
\end{minipage} & \begin{minipage}[b]{\linewidth}\raggedright
Observation
\end{minipage} & \begin{minipage}[b]{\linewidth}\raggedright
Defeats
\end{minipage} \\
\midrule\noalign{}
\endhead
\bottomrule\noalign{}
\endlastfoot
RCE-F1 & On the held-out slice, G0 to G2 with frozen rules resolve under
10\% of tasks with 0 model calls & The claim that most fits already
exist for repair at this scale \\
RCE-F2 & Byte-identical or equivalent patches are under half of the
resolved held-out tasks & The reading of the gate as more than
plausibility on the covered tasks \\
RCE-F3 & Under the wall meter, gross joules per solved task of the RCE
solver exceed those of the cheapest model baseline on the tasks both
solve & The second path's joule advantage on covered tasks \\
RCE-F4 & In the convergence protocol of Section 8.6, fewer than half of
the specifications without a stuck input reach a round with zero model
calls within 20 rounds & Compiling the model out of the loop \\
\end{longtable}
}

\subsection{11. Work to finish}\label{work-to-finish}

Each item is open work with its protocol named.

\begin{enumerate}
\def\labelenumi{\arabic{enumi}.}
\tightlist
\item
  Held-out results on L2 to L4, the Defects4J set of 288 bugs (Section
  8.1), and on L5 to L7 (Section 8.2).
\item
  Baselines in one table, cost-of-pass in joules and dollars (Section
  8.5).
\item
  The first \texttt{measured\_j}: a wall meter on an idle host, at least
  10 repeats, gross and net side by side (Section 8.3).
\item
  The full-file L1 Java result, with a type-indexed API index as the
  next gear (Section 6.2).
\item
  The 626-language survey published with its method (Figure 1).
\item
  New held-out corpora with the gears frozen: Multi-SWE-bench and
  BugsInPy (Section 8.4).
\item
  Convergence of model calls across at least 1,000 specifications
  (Section 8.6).
\end{enumerate}

\subsection{12. Conclusion}\label{conclusion}

A program is a database functional state machine, and RCE treats coding
as a conversion from high-entropy intent to that low-entropy Form.
Correctness is the gate and joules are the ranking. On held-out test
splits with a frozen engine and no model call, the gears repair 99.95\%
of 520,821 one-token faults and all 23,015 in repositories never tuned
on, 38.6\% of 24,072 single-statement Java bugs and 17.4\% of 157,089
Python bugs, at hundredths of a reported joule per solved task on a
package counter. The verification layer proves the rewrites in Z3 and
the kernel in Kani, and differential testing ties five targets to one
trace. What remains is the upper rungs, convergence at large N,
baselines in the same table, and a wall meter. Those protocols are fixed
in Sections 8.3 to 8.6 and the falsifiers in Section 10.

\subsection{References}\label{references}

{[}1{]} C. E. Shannon, ``Prediction and Entropy of Printed English,''
\emph{Bell System Technical Journal} 30(1), 1951.
\href{https://doi.org/10.1002/j.1538-7305.1951.tb01366.x}{doi:10.1002/j.1538-7305.1951.tb01366.x}

{[}2{]} Y. Gurevich, ``Sequential Abstract State Machines Capture
Sequential Algorithms,'' \emph{ACM Transactions on Computational Logic}
1(1), 2000.
\href{https://doi.org/10.1145/343369.343384}{doi:10.1145/343369.343384}

{[}3{]} A. Blass and Y. Gurevich, ``Abstract State Machines Capture
Parallel Algorithms,'' \emph{ACM Transactions on Computational Logic},
2003.
\href{https://doi.org/10.1145/937555.937561}{doi:10.1145/937555.937561}

{[}4{]} N. Dershowitz and Y. Gurevich, ``A Natural Axiomatization of
Computability and Proof of Church's Thesis,'' \emph{Bulletin of Symbolic
Logic}, 2008.
\href{https://doi.org/10.2178/bsl/1231081370}{doi:10.2178/bsl/1231081370}

{[}5{]} E. F. Codd, ``A Relational Model of Data for Large Shared Data
Banks,'' \emph{Communications of the ACM}, 1970.
\href{https://doi.org/10.1145/362384.362685}{doi:10.1145/362384.362685}

{[}6{]} G. H. Mealy, ``A Method for Synthesizing Sequential Circuits,''
\emph{Bell System Technical Journal}, 1955.
\href{https://doi.org/10.1002/j.1538-7305.1955.tb03788.x}{doi:10.1002/j.1538-7305.1955.tb03788.x}

{[}7{]} E. F. Moore, ``Gedanken-Experiments on Sequential Machines,''
\emph{Automata Studies}, 1956.
\href{https://doi.org/10.1515/9781400882618-006}{doi:10.1515/9781400882618-006}

{[}8{]} K. Kang et al., ``Feature-Oriented Domain Analysis (FODA)
Feasibility Study,'' CMU/SEI-90-TR-021, 1990.

{[}9{]} R. Alur et al., ``Syntax-Guided Synthesis,'' FMCAD 2013.
\href{https://doi.org/10.1109/FMCAD.2013.6679385}{doi:10.1109/FMCAD.2013.6679385}

{[}10{]} D. Harel, ``Statecharts: A Visual Formalism for Complex
Systems,'' \emph{Science of Computer Programming}, 1987.
\href{https://doi.org/10.1016/0167-6423(87)90035-9}{doi:10.1016/0167-6423(87)90035-9}

{[}11{]} R. Tate, M. Stepp, Z. Tatlock and S. Lerner, ``Equality
Saturation,'' POPL 2009.
\href{https://doi.org/10.1145/1480881.1480915}{doi:10.1145/1480881.1480915}

{[}12{]} M. Willsey et al., ``egg: Fast and Extensible Equality
Saturation,'' POPL 2021.
\href{https://doi.org/10.1145/3434304}{doi:10.1145/3434304}

{[}13{]} B. Moseley and P. Marks, ``Out of the Tar Pit,'' 2006.

{[}14{]} M. Fowler, ``Event Sourcing,'' 2005.
https://martinfowler.com/eaaDev/EventSourcing.html

{[}15{]} P. Alvaro et al., ``Dedalus: Datalog in Time and Space,'' LNCS
2011.
\href{https://doi.org/10.1007/978-3-642-24206-9_16}{doi:10.1007/978-3-642-24206-9\_16}

{[}16{]} A. M. Turing, ``On Computable Numbers,'' \emph{Proceedings of
the London Mathematical Society}, 1937.
\href{https://doi.org/10.1112/plms/s2-42.1.230}{doi:10.1112/plms/s2-42.1.230}

{[}17{]} A. Church, ``An Unsolvable Problem of Elementary Number
Theory,'' \emph{American Journal of Mathematics}, 1936.
\href{https://doi.org/10.2307/2371045}{doi:10.2307/2371045}

{[}18{]} P. Wadler, ``Propositions as Types,'' \emph{Communications of
the ACM}, 2015.
\href{https://doi.org/10.1145/2699407}{doi:10.1145/2699407}

{[}19{]} J. M. Hellerstein and P. Alvaro, ``Keeping CALM,''
\emph{Communications of the ACM}, 2020.
\href{https://doi.org/10.1145/3369736}{doi:10.1145/3369736}

{[}20{]} G. Kahn, ``The Semantics of a Simple Language for Parallel
Programming,'' IFIP 1974, pp.~471 to 475.

{[}21{]} H. Husain et al., ``CodeSearchNet Challenge,'' 2019.
arXiv:1909.09436

{[}22{]} D. Kim, J. Nam, J. Song and S. Kim, ``Automatic Patch
Generation Learned from Human-Written Patches'' (PAR), ICSE 2013.
\href{https://doi.org/10.1109/ICSE.2013.6606626}{doi:10.1109/ICSE.2013.6606626}

{[}23{]} K. Liu et al., ``TBar: Revisiting Template-based Automated
Program Repair,'' ISSTA 2019. arXiv:1903.08409

{[}24{]} K. Liu et al., ``AVATAR: Fixing Semantic Bugs with Fix Patterns
of Static Analysis Violations,'' SANER 2019. arXiv:1812.07270

{[}25{]} A. Koyuncu et al., ``FixMiner: Mining Relevant Fix Patterns for
Automated Program Repair,'' 2018. arXiv:1810.01791

{[}26{]} Y. Padioleau et al., ``Documenting and Automating Collateral
Evolutions in Linux Device Drivers'' (Coccinelle), EuroSys 2008.
\href{https://doi.org/10.1145/1352592.1352618}{doi:10.1145/1352592.1352618}

{[}27{]} A. Solar-Lezama et al., ``Combinatorial Sketching for Finite
Programs,'' ASPLOS 2006.
\href{https://doi.org/10.1145/1168857.1168907}{doi:10.1145/1168857.1168907}

{[}28{]} L. Chen, M. Zaharia and J. Zou, ``FrugalGPT,'' 2023.
arXiv:2305.05176

{[}29{]} S. Haber and W. S. Stornetta, ``How to Time-Stamp a Digital
Document,'' \emph{Journal of Cryptology}, 1991.
\href{https://doi.org/10.1007/BF00196791}{doi:10.1007/BF00196791}

{[}30{]} S. A. Crosby and D. S. Wallach, ``Efficient Data Structures for
Tamper-Evident Logging,'' USENIX Security 2009.

{[}31{]} L. de Moura and N. Bjørner, ``Z3: An Efficient SMT Solver,''
TACAS 2008.
\href{https://doi.org/10.1007/978-3-540-78800-3_24}{doi:10.1007/978-3-540-78800-3\_24}

{[}32{]} Kani Rust Verifier. https://github.com/model-checking/kani

{[}33{]} C. V. Lopes et al., ``DéjàVu: A Map of Code Duplicates on
GitHub,'' OOPSLA 2017.
\href{https://doi.org/10.1145/3133908}{doi:10.1145/3133908}

{[}34{]} M. Gabel and Z. Su, ``A Study of the Uniqueness of Source
Code,'' FSE 2010.
\href{https://doi.org/10.1145/1882291.1882315}{doi:10.1145/1882291.1882315}

{[}35{]} A. Hindle et al., ``On the Naturalness of Software,'' ICSE
2012.
\href{https://doi.org/10.1109/ICSE.2012.6227135}{doi:10.1109/ICSE.2012.6227135}

{[}36{]} D. Kocetkov et al., ``The Stack: 3 TB of Permissively Licensed
Source Code,'' 2022. arXiv:2211.15533; The Stack v2 dataset card,
https://huggingface.co/datasets/bigcode/the-stack-v2

{[}37{]} R. Just, D. Jalali and M. D. Ernst, ``Defects4J: A Database of
Existing Faults to Enable Controlled Testing Studies for Java
Programs,'' ISSTA 2014.
\href{https://doi.org/10.1145/2610384.2628055}{doi:10.1145/2610384.2628055}

{[}38{]} C. E. Jimenez et al., ``SWE-bench: Can Language Models Resolve
Real-World GitHub Issues?'' 2023. arXiv:2310.06770

{[}39{]} R. Aleithan et al., ``SWE-Bench+: Enhanced Coding Benchmark for
LLMs,'' 2024. arXiv:2410.06992

{[}40{]} C. E. Shannon, ``A Mathematical Theory of Communication,''
\emph{Bell System Technical Journal} 27, 1948.
\href{https://doi.org/10.1002/j.1538-7305.1948.tb01338.x}{doi:10.1002/j.1538-7305.1948.tb01338.x}

{[}41{]} A. N. Kolmogorov, ``Three approaches to the quantitative
definition of information,'' \emph{Problemy Peredachi Informatsii} 1(1),
1965 (in Russian); English translation, \emph{International Journal of
Computer Mathematics} 2, 1968.
\href{https://doi.org/10.1080/00207166808803030}{doi:10.1080/00207166808803030}

{[}42{]} R. J. Solomonoff, ``A formal theory of inductive inference.
Part I,'' \emph{Information and Control} 7(1), 1964.
\href{https://doi.org/10.1016/S0019-9958(64)90223-2}{doi:10.1016/S0019-9958(64)90223-2}

{[}43{]} G. J. Chaitin, ``A Theory of Program Size Formally Identical to
Information Theory,'' \emph{Journal of the ACM} 22(3), 1975.
\href{https://doi.org/10.1145/321892.321894}{doi:10.1145/321892.321894}

{[}44{]} J. Rissanen, ``Modeling by shortest data description,''
\emph{Automatica} 14(5), 1978.
\href{https://doi.org/10.1016/0005-1098(78)90005-5}{doi:10.1016/0005-1098(78)90005-5}

{[}45{]} B. Ray, V. Hellendoorn, S. Godhane, Z. Tu, A. Bacchelli and P.
Devanbu, ``On the `naturalness' of buggy code,'' ICSE 2016.
\href{https://doi.org/10.1145/2884781.2884848}{doi:10.1145/2884781.2884848}

{[}46{]} Z. Tu, Z. Su and P. Devanbu, ``On the localness of software,''
FSE 2014.
\href{https://doi.org/10.1145/2635868.2635875}{doi:10.1145/2635868.2635875}

{[}47{]} M. Allamanis, E. T. Barr, P. Devanbu and C. Sutton, ``A Survey
of Machine Learning for Big Code and Naturalness,'' \emph{ACM Computing
Surveys} 51(4), 2018.
\href{https://doi.org/10.1145/3212695}{doi:10.1145/3212695}

{[}48{]} 陈浙哲, 刘忠鑫, 鄢萌, 徐洲, 夏鑫, 雷晏 (Z. Chen, Z. Liu, M.
Yan, Z. Xu, X. Xia and Y. Lei), ``代码自然性及其应用研究进展'' (Research
progress of code naturalness and its application), \emph{软件学报
(Journal of Software)} 33(8), 3015 to 3034, 2022.
\href{https://doi.org/10.13328/j.cnki.jos.006355}{doi:10.13328/j.cnki.jos.006355}

{[}49{]} G. Sridhara, V. S. Sinha and S. Mani, ``Naturalness of Natural
Language Artifacts in Software,'' ISEC 2015 (India Software Engineering
Conference).
\href{https://doi.org/10.1145/2723742.2723758}{doi:10.1145/2723742.2723758}

{[}50{]} P. Devanbu, ``On the naturalness of software,'' keynote, ISEC
2013 (India Software Engineering Conference).
\href{https://doi.org/10.1145/2442754.2442763}{doi:10.1145/2442754.2442763}

{[}51{]} A. Z. H. Yang, S. Kolak, V. J. Hellendoorn, R. Martins et al.,
``Revisiting Unnaturalness for Automated Program Repair in the Era of
Large Language Models,'' ICSE 2025.
\href{https://doi.org/10.1109/ICSE55347.2025.00089}{doi:10.1109/ICSE55347.2025.00089}

{[}52{]} E. W. Dijkstra, ``Guarded commands, nondeterminacy and formal
derivation of programs,'' \emph{Communications of the ACM} 18(8), 1975.
\href{https://doi.org/10.1145/360933.360975}{doi:10.1145/360933.360975}

{[}53{]} N. Wirth, ``Program development by stepwise refinement,''
\emph{Communications of the ACM} 14(4), 1971.
\href{https://doi.org/10.1145/362575.362577}{doi:10.1145/362575.362577}

{[}54{]} C. A. R. Hoare, ``An axiomatic basis for computer
programming,'' \emph{Communications of the ACM} 12(10), 1969.
\href{https://doi.org/10.1145/363235.363259}{doi:10.1145/363235.363259}

{[}55{]} J.-R. Abrial, \emph{The B-Book: Assigning Programs to
Meanings}, Cambridge University Press, 1996.
\href{https://doi.org/10.1017/CBO9780511624162}{doi:10.1017/CBO9780511624162}

{[}56{]} J.-R. Abrial, \emph{Modeling in Event-B: System and Software
Engineering}, Cambridge University Press, 2010.
\href{https://doi.org/10.1017/CBO9781139195881}{doi:10.1017/CBO9781139195881}

{[}57{]} R.-J. Back and J. von Wright, \emph{Refinement Calculus: A
Systematic Introduction}, Springer, 1998.
\href{https://doi.org/10.1007/978-1-4612-1674-2}{doi:10.1007/978-1-4612-1674-2}

{[}58{]} F. L. Bauer and H. Wössner, \emph{Algorithmic Language and
Program Development}, Springer, 1982.
\href{https://doi.org/10.1007/978-3-642-61807-9}{doi:10.1007/978-3-642-61807-9}

{[}59{]} S. Gulwani, ``Automating string processing in spreadsheets
using input-output examples,'' POPL 2011.
\href{https://doi.org/10.1145/1926385.1926423}{doi:10.1145/1926385.1926423}

{[}60{]} S. Gulwani, O. Polozov and R. Singh, ``Program Synthesis,''
\emph{Foundations and Trends in Programming Languages} 4(1-2), 2017.
\href{https://doi.org/10.1561/2500000010}{doi:10.1561/2500000010}

{[}61{]} F. P. Brooks, ``No Silver Bullet: Essence and Accidents of
Software Engineering,'' \emph{Computer} 20(4), 1987.
\href{https://doi.org/10.1109/MC.1987.1663532}{doi:10.1109/MC.1987.1663532}

{[}62{]} P. Naur, ``Programming as theory building,''
\emph{Microprocessing and Microprogramming} 15(5), 1985.
\href{https://doi.org/10.1016/0165-6074(85)90032-8}{doi:10.1016/0165-6074(85)90032-8}

{[}63{]} M. M. Lehman, ``Programs, life cycles, and laws of software
evolution,'' \emph{Proceedings of the IEEE} 68(9), 1980.
\href{https://doi.org/10.1109/PROC.1980.11805}{doi:10.1109/PROC.1980.11805}

{[}64{]} P. Zave and M. Jackson, ``Four dark corners of requirements
engineering,'' \emph{ACM Transactions on Software Engineering and
Methodology} 6(1), 1997.
\href{https://doi.org/10.1145/237432.237434}{doi:10.1145/237432.237434}

{[}65{]} R. Landauer, ``Irreversibility and Heat Generation in the
Computing Process,'' \emph{IBM Journal of Research and Development}
5(3), 1961.
\href{https://doi.org/10.1147/rd.53.0183}{doi:10.1147/rd.53.0183}

{[}66{]} C. H. Bennett, ``Logical Reversibility of Computation,''
\emph{IBM Journal of Research and Development} 17(6), 1973.
\href{https://doi.org/10.1147/rd.176.0525}{doi:10.1147/rd.176.0525}

{[}67{]} M. Franz and T. Kistler, ``Slim binaries,''
\emph{Communications of the ACM} 40(12), 1997.
\href{https://doi.org/10.1145/265563.265576}{doi:10.1145/265563.265576}

{[}68{]} C. W. Fraser, ``Automatic inference of models for statistical
code compression,'' PLDI 1999.
\href{https://doi.org/10.1145/301618.301672}{doi:10.1145/301618.301672}

{[}69{]} P. Cousot and R. Cousot, ``Abstract interpretation: a unified
lattice model for static analysis of programs by construction or
approximation of fixpoints,'' POPL 1977.
\href{https://doi.org/10.1145/512950.512973}{doi:10.1145/512950.512973}

{[}70{]} Y. Futamura, ``Partial evaluation of computation process: an
approach to a compiler-compiler,'' \emph{Systems, Computers, Controls}
2(5), 1971; reprinted in \emph{Higher-Order and Symbolic Computation}
12(4), 1999.
\href{https://doi.org/10.1023/A:1010095604496}{doi:10.1023/A:1010095604496}

{[}71{]} A. P. Ershov, ``Mixed computation: potential applications and
problems for study,'' \emph{Theoretical Computer Science} 18(1), 1982.
\href{https://doi.org/10.1016/0304-3975(82)90111-6}{doi:10.1016/0304-3975(82)90111-6}

{[}72{]} W. A. Howard, ``The formulae-as-types notion of construction,''
in J. P. Seldin and J. R. Hindley (eds.), \emph{To H. B. Curry: Essays
on Combinatory Logic, Lambda Calculus and Formalism}, Academic Press,
1980.

{[}73{]} 土居真之, 肥後芳樹, 有馬諒, 下仲健斗, 楠本真二 (M. Doi, Y.
Higo, R. Arima, K. Shimonaka and S. Kusumoto),
``ソースコードの「自然さ」を利用した自動生成ファイルの特定''
(Identifying auto-generated files using the naturalness of source code),
電子情報通信学会技術研究報告 (IEICE Technical Report) 117(380), 19 to
24, January 2018. https://sdl.ist.osaka-u.ac.jp/pman/pman3.cgi?D=517

{[}74{]} M. Doi, Y. Higo, R. Arima, K. Shimonaka and S. Kusumoto, ``On
the naturalness of auto-generated code: can we identify auto-generated
code automatically?'' ICPC 2018.
\href{https://doi.org/10.1145/3196321.3196356}{doi:10.1145/3196321.3196356}

{[}75{]} F. V. Arrebola, ``Modelo de sugestão de código-fonte baseado em
modelos de linguagem organizados em um contexto hierárquico'' (A source
code suggestion model based on language models organized in a
hierarchical context), doctoral thesis in electrical engineering,
Fundação Educacional Inaciana Padre Sabóia de Medeiros (FEI), Brazil,
2018. Advisor R. A. C. Bianchi.

{[}76{]} G. Febres, K. Jaffé and C. Gershenson, ``Complexity measurement
of natural and artificial languages,'' \emph{Complexity} 20(6), 2015.
\href{https://doi.org/10.1002/cplx.21529}{doi:10.1002/cplx.21529}

{[}77{]} M. Rahman, D. Palani and P. C. Rigby, ``Natural Software
Revisited,'' ICSE 2019.
\href{https://doi.org/10.1109/ICSE.2019.00022}{doi:10.1109/ICSE.2019.00022}

{[}78{]} A. Torres, S. Baltes, C. Treude and M. Wagner, ``Applying
Information Theory to Software Evolution,'' NLBSE 2023.
\href{https://doi.org/10.1109/NLBSE59153.2023.00017}{doi:10.1109/NLBSE59153.2023.00017}

{[}79{]} A. Blass, N. Dershowitz and Y. Gurevich, ``When are two
algorithms the same?,'' \emph{Bulletin of Symbolic Logic} 15(2), 2009.
arXiv:0811.0811. https://arxiv.org/abs/0811.0811

{[}80{]} Y. N. Moschovakis, ``What is an algorithm?,'' in
\emph{Mathematics Unlimited: 2001 and Beyond}, Springer, 2001.
https://www.math.ucla.edu/\textasciitilde ynm/papers/eng.pdf

{[}81{]} M. Y. Vardi, ``What is an algorithm?,'' \emph{Communications of
the ACM} 55(3), 2012.
\href{https://doi.org/10.1145/2093548.2093549}{doi:10.1145/2093548.2093549}

{[}82{]} S. Laddad, A. Cheung, J. M. Hellerstein and M. Milano, ``Flo: a
Semantic Foundation for Progressive Stream Processing,'' POPL 2025.
arXiv:2411.08274. https://arxiv.org/abs/2411.08274

{[}83{]} The Elm Architecture, Elm guide.
https://guide.elm-lang.org/architecture/

{[}84{]} Temporal, ``Workflows'' (event history), documentation.
https://docs.temporal.io/workflows

{[}85{]} A. Pnueli, M. Siegel and E. Singerman, ``Translation
validation,'' TACAS 1998.
\href{https://doi.org/10.1007/BFb0054170}{doi:10.1007/BFb0054170}

{[}86{]} X. Yang, Y. Chen, E. Eide and J. Regehr, ``Finding and
understanding bugs in C compilers,'' PLDI 2011.
\href{https://doi.org/10.1145/1993498.1993532}{doi:10.1145/1993498.1993532}

{[}87{]} C. Nandi, M. Willsey, A. Zhu et al., ``Rewrite Rule Inference
Using Equality Saturation,'' OOPSLA 2021. arXiv:2108.10436.
https://arxiv.org/abs/2108.10436

{[}88{]} C. Le Goues, T. Nguyen, S. Forrest and W. Weimer, ``GenProg: A
Generic Method for Automatic Software Repair,'' \emph{IEEE TSE} 38(1),
2012.
\href{https://doi.org/10.1109/TSE.2011.104}{doi:10.1109/TSE.2011.104}

{[}89{]} F. Long and M. Rinard, ``Automatic patch generation by learning
correct code,'' POPL 2016.
\href{https://doi.org/10.1145/2837614.2837617}{doi:10.1145/2837614.2837617}

{[}90{]} J. Bader, A. Scott, M. Pradel and S. Chandra, ``Getafix:
Learning to Fix Bugs Automatically,'' OOPSLA 2019. arXiv:1902.06111.
https://arxiv.org/abs/1902.06111

{[}91{]} K. Etemadi, N. Harrand, S. Larsen et al., ``Sorald: Automatic
Patch Suggestions for SonarQube Static Analysis Violations,'' \emph{IEEE
TDSC}, 2022. arXiv:2103.12033. https://arxiv.org/abs/2103.12033

{[}92{]} Z. Qi, F. Long, S. Achour and M. Rinard, ``An analysis of patch
plausibility and correctness for generate-and-validate patch generation
systems,'' ISSTA 2015.
\href{https://doi.org/10.1145/2771783.2771791}{doi:10.1145/2771783.2771791}

{[}93{]} E. K. Smith, E. T. Barr, C. Le Goues and Y. Brun, ``Is the cure
worse than the disease? Overfitting in automated program repair,''
ESEC/FSE 2015.
\href{https://doi.org/10.1145/2786805.2786825}{doi:10.1145/2786805.2786825}

{[}94{]} X.-B. D. Le, F. Thung, D. Lo and C. Le Goues, ``Overfitting in
semantics-based automated program repair,'' \emph{Empirical Software
Engineering} 23, 2018.
\href{https://doi.org/10.1007/s10664-017-9577-2}{doi:10.1007/s10664-017-9577-2}

{[}95{]} Y. Xiong, X. Liu, M. Zeng, L. Zhang and G. Huang, ``Identifying
patch correctness in test-based program repair,'' ICSE 2018.
\href{https://doi.org/10.1145/3180155.3180182}{doi:10.1145/3180155.3180182}

{[}96{]} T. Durieux, F. Madeiral, M. Martinez and R. Abreu, ``Empirical
Review of Java Program Repair Tools: A Large-Scale Experiment on 2,141
Bugs and 23,551 Repair Attempts,'' ESEC/FSE 2019. arXiv:1905.11973.
https://arxiv.org/abs/1905.11973

{[}97{]} K. Liu, A. Koyuncu, T. F. Bissyandé et al., ``You Cannot Fix
What You Cannot Find! An Investigation of Fault Localization Bias in
Benchmarking Automated Program Repair Systems,'' ICST 2019.
arXiv:1812.07283. https://arxiv.org/abs/1812.07283

{[}98{]} M. Monperrus, ``Automatic Software Repair: A Bibliography,''
\emph{ACM Computing Surveys} 51(1), 2018.
\href{https://doi.org/10.1145/3105906}{doi:10.1145/3105906}

{[}99{]} C. S. Xia, Y. Deng, S. Dunn and L. Zhang, ``Agentless:
Demystifying LLM-based Software Engineering Agents,'' 2024.
arXiv:2407.01489. https://arxiv.org/abs/2407.01489

{[}100{]} S. Kapoor, B. Stroebl, Z. S. Siegel, N. Nadgir and A.
Narayanan, ``AI Agents That Matter,'' 2024. arXiv:2407.01502.
https://arxiv.org/abs/2407.01502

{[}101{]} M. H. Erol, B. El, M. Suzgun, M. Yuksekgonul and J. Zou,
``Cost-of-Pass: An Economic Framework for Evaluating Language Models,''
2025. arXiv:2504.13359. https://arxiv.org/abs/2504.13359

{[}102{]} D. Zan, Z. Huang, W. Liu et al., ``Multi-SWE-bench: A
Multilingual Benchmark for Issue Resolving,'' 2025. arXiv:2504.02605.
https://arxiv.org/abs/2504.02605

{[}103{]} R. Widyasari, S. Q. Sim, C. Lok et al., ``BugsInPy: a database
of existing bugs in Python programs to enable controlled testing and
debugging studies,'' ESEC/FSE 2020.
\href{https://doi.org/10.1145/3368089.3417943}{doi:10.1145/3368089.3417943}

{[}104{]} M. Allamanis, ``The Adverse Effects of Code Duplication in
Machine Learning Models of Code,'' Onward! 2019. arXiv:1812.06469.
https://arxiv.org/abs/1812.06469

{[}105{]} K. N. Khan, M. Hirki, T. Niemi, J. K. Nurminen and Z. Ou,
``RAPL in Action: Experiences in Using RAPL for Power Measurements,''
\emph{ACM TOMPECS} 3(2), 2018.
\href{https://doi.org/10.1145/3177754}{doi:10.1145/3177754}

{[}106{]} The Rust project, rustfix and \texttt{cargo\ fix}.
https://github.com/rust-lang/rustfix

{[}107{]} R.-M. Karampatsis and C. Sutton, ``How Often Do
Single-Statement Bugs Occur? The ManySStuBs4J Dataset,'' MSR 2020.
\href{https://doi.org/10.1145/3379597.3387491}{doi:10.1145/3379597.3387491}

{[}108{]} C. Richter and H. Wehrheim, ``TSSB-3M: Mining single statement
bugs at massive scale,'' MSR 2022.
\href{https://doi.org/10.1145/3524842.3528505}{doi:10.1145/3524842.3528505}

{[}109{]} H. Coles, T. Laurent, C. Henard, M. Papadakis and A.
Ventresque, ``PIT: a practical mutation testing tool for Java,'' ISSTA
2016.
\href{https://doi.org/10.1145/2931037.2948707}{doi:10.1145/2931037.2948707}

{[}110{]} R. Just, D. Jalali, L. Inozemtseva, M. D. Ernst, R. Holmes and
G. Fraser, ``Are mutants a valid substitute for real faults in software
testing?'' FSE 2014.
\href{https://doi.org/10.1145/2635868.2635929}{doi:10.1145/2635868.2635929}

\end{document}
